% \section{Trustworthiness and Reputation Models in Crowdsourced Systems}

% A fundamental challenge in blockchain-based civic reporting systems is enabling anonymous participation while maintaining trust in submitted reports. Traditional reporting systems derive trust from user identity; however, decentralized and privacy-preserving systems must instead evaluate trust based on user behavior and collective validation. The literature addresses this challenge through reputation and trustworthiness models that decouple identity from reliability, allowing users to remain anonymous while their contributions are evaluated objectively.

% \subsubsection{Concept of Trustworthiness in Crowdsourced Reporting}

% Trustworthiness in crowdsourced reporting is commonly defined as a multi-dimensional concept encompassing data quality, consistency, and contributor behavior \cite{gheyas2025establishing}. Data quality refers to the accuracy and completeness of submitted information, data reliability is assessed through agreement with other reports or historical patterns, and contributor credibility reflects the past behavior of the reporting entity. Together, these dimensions allow systems to assess reports without requiring real-world identity disclosure.

% \subsubsection{Anonymous Reputation Mechanisms}

% To support anonymity while preventing abuse, several reputation mechanisms have been proposed. Identity-bound pseudonymity is a widely cited approach, where a user possesses a verified root identity but interacts with the system using unlinkable pseudonyms \cite{androulaki2008reputation}. Reputation is cryptographically linked to the root identity, ensuring accountability while preserving privacy. This approach mitigates Sybil attacks, as malicious behavior under one pseudonym affects the user's overall reputation, limiting repeated abuse \cite{androulaki2008reputation}, \cite{kaafarani2018anonymous}.

% \subsubsection{Blockchain-Based Reputation Frameworks}

% The PriRPT framework is a notable example of a privacy-preserving reporting system that integrates anonymous credentials and reputation mechanisms \cite{lin2023prirpt}. Users register with a trusted authority and subsequently submit reports using derived anonymous credentials. Verified reports are rewarded without linking the reporter’s identity to the reward, preserving unlinkability. By employing optimized cryptographic schemes, PriRPT demonstrates feasibility for mobile-based civic reporting applications.

% Similarly, the Block-CITE framework proposes a dual-blockchain architecture to separate reporting tasks from reputation management \cite{li2025blockcite}. One chain stores report data, while a second chain maintains reputation scores. This separation improves scalability and allows reputation updates without congesting the primary reporting ledger. Reputation is incremented based on report validation outcomes, reinforcing honest participation while maintaining privacy.

% \subsubsection{Peer Validation and Consensus-Based Trust}

% In decentralized environments, trust is often established through peer validation mechanisms. Reports may be verified by nearby users or relevant stakeholders, leveraging the “wisdom of the crowd” principle \cite{hasan2023privacy}. Consensus-based approaches reward users whose validations align with the majority, while penalizing inconsistent behavior. To prevent manipulation, many systems employ weighted voting schemes where votes are scaled by reputation scores, ensuring that trusted contributors have greater influence \cite{leonardos2022weighted}, \cite{aluko2021proof}, \cite{hewa2025repuchain}.

% \subsubsection{Algorithmic Trust and Reputation Scoring}

% Trust scores are typically computed using mathematical models that adapt over time. Subjective logic models represent trust as a combination of belief, disbelief, and uncertainty, allowing nuanced decision-making in uncertain environments \cite{zupancic2019data} ,\cite{fu2022bfcri}. Other approaches use growth functions and time-decay mechanisms to ensure that reputation is earned gradually and must be maintained through consistent honest behavior \cite{zupancic2019data}, \cite{das2014re3}. These models reduce the risk of reputation manipulation and ensure responsiveness to recent behavior.

% \subsubsection{Trust in Mobile Crowdsensing Contexts}

% For automated or sensor-based reports, such as those generated by mobile devices, trust evaluation must account for device reliability and environmental variability. Probabilistic models compare sensor readings against nearby observations to detect anomalies or faulty data \cite{zupancic2019data}, \cite{ali2020blockchain}. Additionally, location verification techniques, including proof-of-location protocols, are used to ensure that reports originate from the claimed physical location, preventing fraudulent submissions \cite{albilali2024blockchain}.

% Overall, the literature demonstrates that combining anonymous reputation systems, peer validation, and adaptive trust models enables reliable crowdsourced reporting without compromising user privacy. These approaches directly inform the design of blockchain-based community reporting platforms that aim to balance anonymity, accountability, and data integrity.


% % \subsection{Trustworthiness and Reputation Models in Crowdsourced Systems}

% % The core value proposition of a blockchain-based reporting system is that it allows for anonymous yet trusted participation. This presents a seeming paradox: how can a system trust a report if the source is unknown? In traditional systems, trust is transitive through identity (I trust User X, therefore I trust their report). In anonymous decentralized systems, trust must be derived from behavior and consensus. The literature resolves this through Reputation Systems and Trustworthiness Evaluation Models that dissociate identity (who you are) from reputation (how reliable you have been).

% % \subsubsection{Defining Trustworthiness in Crowdsourcing}

% % Trustworthiness in this context is a composite metric comprising three dimensions \cite{gheyas2025establishing}:  

% % \begin{enumerate}
% %     \item Data Quality: Is the data accurate, complete, and formatted correctly? (Objective verification).
% %     \item Data Reliability: Is the data consistent with other observations and historical trends? (Consensus verification). 
% %     \item Contributor Credibility: Does the user have a history of honest contributions? (Reputation).  
% % \end{enumerate}

% % \subsubsection{Mechanisms for Anonymous Reputation}

% % To evaluate trustworthiness without revealing identity, researchers have developed sophisticated cryptographic and algorithmic frameworks that allow a user to prove they are "reputable" without proving "who" they are.

% % \subsubsection{ Identity-Bound Pseudonymity}

% % A key innovation is the "Identity-Bound Reputation System" \cite{androulaki2008reputation}. In this model, a user has a real identity (e.g., a master key registered with a government authority) but interacts with the system using disposable pseudonyms. The reputation is mathematically linked to the master identity, but the link is hidden from the public .

% % \begin{itemize}
% %     \item Mechanism: When a user submits a report, they generate a Zero-Knowledge Proof (ZKP) that attests: "I am a member of the set of registered citizens, and my reputation score is above Threshold X," without revealing which member they are.
% %     \item Unlinkability: This prevents "Sybil attacks" (one user creating multiple fake accounts to boost a rating or spam the system) because the reputation is tied to a singular, verifiable root identity. If a user acts maliciously under one pseudonym, the penalty propagates to their master reputation, affecting all future pseudonyms \cite{androulaki2008reputation} \cite{kaafarani2018anonymous}. 
% % \end{itemize}

% % \subsubsection{The "PriRPT" Framework}

% % The PriRPT (Practical blockchain-based Privacy-Preserving Reporting System) is a prominent example in the literature that specifically addresses this challenge \cite{lin2023prirpt}. It integrates Keyed-Verification Anonymous Credentials (KVAC) and Structure-Preserving Signatures on Equivalence Classes (SPS-EQ).

% % Workflow:

% % \begin{enumerate}
% %     \item Registration: User registers with a trusted authority (TA) and receives a credential.
% %     \item Reporting: User submits a report signed with a derived credential. The signature proves validity but not identity.
% %     \item Rewarding: If the report is verified (e.g., leads to a fix), the system issues a reward token. The user can claim this token anonymously using the SPS-EQ scheme, breaking the link between the report submission and the reward payout address.
% %     \item Efficiency: By replacing general-purpose ZKPs with optimized signature schemes (SPS-EQ), PriRPT reduces the computational overhead, making it feasible for mobile devices, which is a critical constraint for citizen reporting apps \cite{lin2023prirpt}. 
% % \end{enumerate}

% % \subsubsection{The "Block-CITE" Framework}

% % Block-CITE proposes a dual-blockchain architecture to separate the task (the report) from the reputation (the trust score) to enhance scalability and privacy \cite{li2025blockcite}. 

% % \begin{itemize}
% %     \item TaskChain: Records the details of the crowdsourcing task (e.g., the report, evidence, and timestamps).
% %     \item RepuChain: Records the reputation scores of the participants.
% %     \item Interaction: Reputation is updated based on the outcome of the task. If a report is validated by the government or peers, the reputation score in RepuChain is incremented via a smart contract. This separation ensures that frequent reputation updates do not clutter the TaskChain, and allows for different privacy settings on each chain \cite{li2025blockcite}.
% % \end{itemize}

% % \subsubsection{Peer Validation and "Wisdom of the Crowd"}

% % In the absence of a central validator, the system must rely on peer validation. This is often modeled using Game Theory and Weighted Voting.

% % \begin{itemize}
% %     \item Schelling Points: In a decentralized voting game, users are incentivized to vote with the majority. If a user reports a pothole, and 10 neighbors also verify it (consensus), the report is deemed true. Users who vote with the consensus are rewarded; those who vote against are penalized. This mechanism assumes that the majority of citizens are honest \cite{hasan2023privacy}.  
% %     \item Weighted Voting: To prevent "mob rule" or attacks by bot farms, votes are weighted by reputation. A vote from a user with a high "Trustworthiness Score" counts significantly more than a new user's vote. This is detailed in "Proof of Reputation" (PoR) consensus mechanisms \cite{leonardos2022weighted}. The RepuChain protocol, for instance, uses a multi-factor reputation score (incorporating data validity, frequency, and latency) to elect leaders in the consensus process, ensuring that only trusted nodes validate data \cite{aluko2021proof} \cite{hewa2025repuchain}. 
% % \end{itemize}

% % \subsubsection{Algorithmic Trust Models}

% % Mathematical models are used to calculate the trust score dynamically, ensuring it reflects current behavior.

% % \begin{itemize}
% %     \item Subjective Logic: This approach models trust not as a single scalar number but as a "belief, disbelief, and uncertainty" triplet. It accounts for the fact that trust is subjective and transitive (User A might trust User B, but User C might not). This allows the system to filter reports based on a user's personal trust network \cite{zupancic2019data} \cite{fu2022bfcri}.   
% %     \item Gompertz Function: Some systems use the Gompertz function (an S-shaped curve) to model reputation growth—reputation is easy to lose (steep drop) but hard to gain (slow climb). This prevents users from rapidly "gaming" the system to achieve high trust scores before launching an attack \cite{zupancic2019data}. 
% %     \item Time-Decay: Reputation models often include a time-decay factor (EWMA - Exponentially Weighted Moving Average), ensuring that past good behavior doesn't grant permanent immunity to a user who starts acting maliciously. A user must maintain constant honest participation to keep a high score \cite{das2014re3}.
% % \end{itemize}

% % \subsubsection{Trust in Mobile Crowdsensing (MCS)}

% % For automated reports (e.g., a phone's accelerometer detecting a pothole automatically), trust models must account for device calibration and environmental noise.

% % \begin{itemize}
% %     \item Probabilistic Assessment: Trust is inferred from the statistical deviation of a user's data from the group mean. If one sensor reports a temperature of 100°C while ten others nearby report 25°C, the outlier is flagged as untrustworthy (faulty or malicious) \cite{zupancic2019data} \cite{ali2020blockchain}.   
% %     \item Location Proofs: "Proof of Location" protocols (using GPS + Wi-Fi triangulation logs stored on-chain) are critical to ensure that the user was actually physically present at the site of the report, preventing "armchair reporting" where users spoof locations to file fake reports \cite{albilali2024blockchain}.
% % \end{itemize}

